% Blatt 7, A1. Werte und Aufgabenstellung: aufgaben/Blatt07.pdf, S. 1.

\subsection*{Gegeben, Modell und Vorzeichen}
Die Scheibe hat die Masse $m=\qty{0.50}{\kilogram}$ und den Radius
$R=\qty{0.050}{\metre}$. Der praktisch masselose Stab hat den Radius
$r=\qty{0.0050}{\metre}$, die Schnüre haben die Länge
$l=\qty{0.50}{\metre}$. Wir verwenden $g=\qty{9.81}{\metre\per\second\squared}$.
Der Zahlenwert für $g$ ist eine verwendete Näherung, keine zusätzliche Angabe
des Aufgabenblatts.

Für die Bewegung ausserhalb des Umkehrvorgangs nehmen wir masselose,
undehnbare Schnüre, schlupffreies Ab- und Aufrollen und vernachlässigbare
Reibung an. Die Absenkung des Schwerpunkts bis kurz vor dem Ende des Abrollens
setzen wir gleich der abgewickelten Länge $l$. Die Umkehr selbst betrachten wir
getrennt, denn ihre Dauer und ihr Mechanismus sind nicht angegeben.

$T$ sei die \emph{gesamte} nach oben auf das Rad wirkende Zugkraft beider
Schnüre. Bei symmetrischer Aufhängung trägt jede Schnur $T/2$. Die Kraft des
Rades auf die Stativstange zeigt entgegengesetzt, also nach unten, und hat bei
masselosen Schnüren den Betrag $T$.

\subsection*{a) Beschleunigung beim Abrollen}
Als positive Richtung für die Schwerpunktbewegung wählen wir nach unten und
für die Drehbeschleunigung die Abrollrichtung. Gewicht und Schnurzug wirken auf
den Schwerpunkt entgegengesetzt. Die Schnüre greifen am Stabradius $r$ an und
üben dort ein Drehmoment aus. Damit gelten
\begin{equation}
  mg-T=ma, \qquad Tr=I_S\alpha, \qquad a=r\alpha,
  \qquad I_S=\frac12 mR^2.
\end{equation}
Der letzte Zusammenhang ist das Trägheitsmoment einer homogenen Scheibe um
ihre Symmetrieachse. Die Gewichtskraft selbst erzeugt um den Schwerpunkt kein
Drehmoment. Mit $\alpha=a/r$ ergibt die Drehgleichung $T=I_Sa/r^2$. Einsetzen
in die Schwerpunktgleichung liefert
\begin{align}
  mg&=a\left(m+\frac{I_S}{r^2}\right),\\
  a&=\frac{g}{1+I_S/(mr^2)}
    =\frac{g}{1+R^2/(2r^2)}.
\end{align}
Daraus folgt
\begin{ergebnis}[a) Beschleunigung]
\begin{equation*}
  a=\frac{\qty{9.81}{\metre\per\second\squared}}{51}
   \approx\qty{0.192}{\metre\per\second\squared}\quad\text{abwärts}.
\end{equation*}
\end{ergebnis}
Das Rad fällt langsamer als im freien Fall: Ein Teil der Bewegung ist Rotation.

\subsection*{b) Kraft auf die Stativstange}
Aus der Schwerpunktgleichung folgt $T=m(g-a)$. Die beiden Schnurkräfte
zusammen ziehen die Stange nach unten:
\begin{ergebnis}[b) Gesamtkraft auf die Stange]
\begin{equation*}
  F_{\mathrm{Stange}}=T
    =\qty{0.50}{\kilogram}\,
      (\qty{9.81}{\metre\per\second\squared}
       -\qty{0.19235}{\metre\per\second\squared})
    \approx\qty{4.81}{\newton}\quad\text{abwärts}.
\end{equation*}
\end{ergebnis}
Bei gleicher Lastverteilung sind das ungefähr \qty{2.40}{\newton} pro Schnur.
Die angegebene Gesamtkraft ist kleiner als die Gewichtskraft
$mg=\qty{4.905}{\newton}$, weil das Rad abwärts beschleunigt. Das Eigengewicht
von Stange und Schnüren ist nicht Teil dieser Kraft.

\subsection*{c) Geschwindigkeit kurz vor dem Ende des Abrollens}
Das Rad startet aus der Ruhe. Da $a$ im idealisierten Abrollmodell konstant
ist und der Schwerpunkt die Strecke $l$ zurücklegt, gilt
\begin{equation}
  v^2=2al
  =\frac{2gl}{1+R^2/(2r^2)}.
\end{equation}
\begin{ergebnis}[c) Geschwindigkeit]
\begin{equation*}
  v=\sqrt{2\cdot\qty{0.19235}{\metre\per\second\squared}
              \cdot\qty{0.50}{\metre}}
    \approx\qty{0.439}{\metre\per\second}\quad\text{abwärts}.
\end{equation*}
\end{ergebnis}
Als unabhängige Kontrolle dient die Energiegleichung
$mgl=\tfrac12 mv^2+\tfrac12 I_S(v/r)^2$: Sie berücksichtigt sowohl die
Geschwindigkeit des Schwerpunkts als auch die Drehung der Scheibe.

\subsection*{d) Mittlere zusätzliche Kraft bei der Umkehr}
Hier reichen die Angaben des Blattes für keinen eindeutigen Zahlenwert der
mittleren Kraft. Weder die Dauer des Umkehrvorgangs noch die Geschwindigkeit
danach oder die Energieverluste sind angegeben. Je kürzer die Umkehr bei gleicher
Impulsänderung dauert, desto grösser ist die mittlere Kraft.

Sei $\Delta t$ die Dauer, $v_-$ der Betrag der Geschwindigkeit unmittelbar
vor der Umkehr und $v_+$ der Betrag unmittelbar danach. Mit positiver
Aufwärtsrichtung ist die Impulsänderung des Rades
$\Delta p=m(v_+ + v_-)$. Im Modell, in dem vertikal nur die gesamte
Schnurkraft $T(t)$ und die Gewichtskraft wirken, folgt
\begin{equation}
  (\overline T-mg)\Delta t=m(v_+ + v_-),
  \qquad
  \overline T=mg+\frac{m(v_+ + v_-)}{\Delta t}.
\end{equation}
Die gleich grosse mittlere Zugkraft wirkt nach unten auf die Stativstange.
\emph{Zusätzlich zum statischen Gewicht $mg$} ist sie daher
$m(v_+ + v_-)/\Delta t$. Wenn stattdessen die Last während des Abrollens aus
Teil b als Bezug gemeint ist, beträgt die Zunahme
$ma+m(v_+ + v_-)/\Delta t$.

Unter der \emph{zusätzlichen} idealisierten Annahme einer symmetrischen,
verlustfreien Umkehr mit gleichen Geschwindigkeitsbeträgen unmittelbar davor
und danach wäre $v_+=v_-=v$. Dann wäre der zusätzliche Kraftstoss relativ zum
Gewicht $2mv\approx\qty{0.439}{\newton\second}$. Eine mittlere Zusatzkraft
folgt daraus erst, wenn $\Delta t$ bekannt ist. Der Kraftstoss ist kein
Kraftwert. Im Grenzmodell einer sprungartigen Umkehr ist keine endliche mittlere
Kraft für eine verschwindende Zeitdauer bestimmbar.

\subsection*{e) Bonus: Kraft beim anschliessenden Aufrollen}
Nach dem eigentlichen Umkehrvorgang bewegt sich der Schwerpunkt aufwärts,
während das Rad durch die Schwerkraft abgebremst wird. Für die ideale
Aufrollphase nehmen wir wieder eine gespannte, schlupffrei auf der anderen
Seite der Achse auflaufende Schnur und vernachlässigbare Reibung an. Dann
zeigt die Schwerpunktbeschleunigung erneut nach unten und hat den gleichen
Betrag wie in Teil a. Mit $mg-T_{\mathrm{auf}}=ma$ folgt
\begin{ergebnis}[e) Gesamtkraft beim Aufrollen]
\begin{equation*}
  F_{\mathrm{Stange,auf}}=T_{\mathrm{auf}}=m(g-a)
   \approx\qty{4.81}{\newton}\quad\text{abwärts}.
\end{equation*}
\end{ergebnis}
Das gilt für die reguläre Aufrollphase, \emph{nicht} für die kurzzeitige
Umkehr. In dieser Phase sind Geschwindigkeit und Beschleunigung
entgegengesetzt gerichtet. Verluste, Schlupf oder ein anderer Kontaktmechanismus
können die reale Kraft ändern.

\begin{kontrolle}[Plausibilität und Geltungsbereich]
Das Ergebnis für $a$ liegt zwischen null und $g$; $T$ ist positiv und kleiner
als $mg$. Die Energiegleichung reproduziert die Geschwindigkeit aus der
Kinematik. Die Zahl in Teil d lässt sich ohne Angaben zur Umkehrzeit und zum
Verlauf nicht bestimmen. Die Werte für die übrigen Teilaufgaben sind Aussagen
über das angegebene ideale Modell, nicht über jede reale Apparatur.
\end{kontrolle}

\begin{merksatz}
Bei rollenden oder aufgewickelten Körpern müssen Schwerpunktbewegung und
Rotation zusammenpassen. Bestimme zuerst Kräfte und wirksame Hebelarme; eine
Kraft während eines kurzen Stosses erfordert zusätzlich eine Aussage über die
Dauer oder den zeitlichen Verlauf.
\end{merksatz}
